\chapter{The particle as a persistent material section: history, representation, band, composition, and identity criterion}
\label{chap:particula-seccion-material-persistente}

\section{Definition of particle as a compatible class of persistent material sections under refinement}
A particle is a compatible class of persistent material sections of the genealogical extension, endowed with an observable history, representation, orientation, and a criterion of identity under refinement. The route is its history, not the particle itself; the cell is a coordinate of the atlas, not a moving body; mass is a spectral observable, not its ontological definition.

\section{Classifier of material species by fibre, representation, charge, spin, flavour and generation}
The species is typed by charge, spin, color, flavor, generation,
representation, and orientation. Charge is an oriented character; spin is
representation data; color is a residual ternary fiber; flavor and
generation belong to the multisection and its depth. None of these
data reduces to the mass value.

\section{Antiparticle, Band, and Virtuality}
The involution \(C_M=-\operatorname{rev}\) induces the anti-section and
sheet conjugation. Particle and antiparticle may be two sections of the
same Möbius cover when the corresponding loop, holonomy, and
representation are fixed. The extended return of 216 steps does not
decompose into two independent trajectories of 108. A virtual particle is
an intermediate nonasymptotic section with a finite horizon of prolongation.

\section{Path persistence, multisections, and Möbius conjugation}

Persistence is formulated as a section of an inverse system; the internal
species as an orbital quotient of a multisection; and conjugation as an
anti-section with declared holonomy. The complete development preserves the
extension--route--record--signature stratification, the temporal reader, and the
distinction between the two occurrences of \(2\pi/4\pi\).

\begingroup
\let\chapter\section
\let\section\subsection
\let\subsection\subsubsection
\let\subsubsection\paragraph
\let\clearpage\relax
\let\cleardoublepage\relax
\section{Route persistence, species multisections, and conjugation of Möbius type}
\label{sec:persistencia-ruta-mobius-familias}

This section specifies a distinction that, if formulated too
briefly, reverses the logic of the construction. The route is not a datum that
a particle chooses once its mass is known. Nor is it an auxiliary drawing
without causal content. It is the observable and oriented history of a surviving extension.
The extension additionally preserves compatible depth, the
register, carries, and memory that do not fit into a finite cutoff. In this
precise sense,
\[
 \boxed{\text{an HMT particle is the compatible persistence of an
 enriched route}.}
\]
The phrase and the definition by extensions express the same object at two
levels: the route is its visible history; the extension is the complete tower of
that history with all its compatible prolongations.

\subsection{Persistence as a section of an inverse system}

Let \((\mathcal S_m,\rho_{m+1,m})_{m\geq0}\) be the system of states that
survive to depth \(m\), and let
\[
 x=(x_m)_{m\geq0}\in\Omega_{\rm HMT}^{\rm surv}
 =\varprojlim_m\mathcal S_m.
\]
Route shadows satisfy the compatibility square
\[
 \operatorname{tr}_{m+1,m}\!\left(\operatorname{sh}_{m+1}(x_{m+1})\right)
 =\operatorname{sh}_{m}\!\left(\rho_{m+1,m}(x_{m+1})\right)
 =\operatorname{sh}_{m}(x_m).
\]
Therefore
\[
 \gamma_x=\bigl(\gamma_x^{[m]}\bigr)_{m\geq0},
 \qquad \gamma_x^{[m]}=\operatorname{sh}_m(x_m),
\]
is a compatible history, not a sequence retroactively fitted.

Over each cut \(\gamma_x^{[m]}\) one has the discrete bundle
\(\mathcal B_{\gamma_x^{[m]}}\). A persistent section is a family
\[
 s=(s_m)_{m\geq0},
 \qquad
 s_m\in\Gamma\!\left(\mathcal B_{\gamma_x^{[m]}}\right),
 \qquad
 \operatorname{res}_{m+1,m}(s_{m+1})=s_m.
\]
This equation is the mathematical content of persistence: each finite reading
extends without contradicting the preceding readings. The
particle tuple of \cref{chap:particula} adds occupation,
representation, and propagation, but does not alter this provenance.

\subsection{Internal species as an orbital quotient of the multisection}

In the nonadic atlas, each oriented cell \(c=(r,c')\in\mathcal C_{81}\)
carries a route family
\[
 \mathcal R(c)=
 (n_A,s_C,k_{\Delta_4},\nu_{120},\nu_{270},\tau_{12}).
\]
The image of \(\mathcal R\) contains fifty-six values. Once the
chart and its order are fixed, rank \(\kappa\in\{0,1,2,3,4\}\) within each fibre
produces the injective lift
\[
 c\longmapsto \bigl(\mathcal R(c),\kappa(c)\bigr).
\]
The five-state alphabet is minimal, since two fibers of
cardinality five exist. This selection is internal, prospective,
and blind to masses and particle names.

The complete census supporting this statement is
\[
 81=45_{\rm leptones}+36_{\rm quarks}
   =25_{\ell^-}+20_{\nu}+20_{d\text{-type}}+16_{u\text{-type}}.
\]
The fifty-six route fibers possess the histogram.
\[
 41\times1+10\times2+2\times3+1\times4+2\times5=81.
\]
Therefore, \(56\) counts route signatures, whereas \(81\) counts oriented
cells; memory \(\kappa\) distinguishes cells that one and the same signature does not
separate.

To classify internal species, this map is not inverted. The
quotient is taken.
\[
 q:\mathcal C_{81}\longrightarrow\Sigma_{13},
 \qquad
 q(c)=\bigl(\operatorname{familia}(c),G(c)\bigr),
\]
whose image comprises thirteen types. The output associated with a type \(y\) is the
finite multisection of the atlas
\[
 \boxed{
 \mathfrak S_y^{\rm atlas}=
 \{(\mathcal R(c),\kappa(c)):q(c)=y\}.}
\]
Their cardinals are
\[
 1,8,16;\qquad 2,4,6,8;\qquad 2,4,6,8;\qquad4,12,
\]
for charged leptons, neutrinos, and the two types of quark,
respectively. The species is the common type of that fiber; a particular particle
is a persistent section realized within it. Requiring the
name of a species by itself to select a cell destroys precisely the
degrees of orientation and conjugation that the fiber preserves.

The elevation to extensions is a different typed object. If
\(\pi_{81}:\Omega_{\rm HMT}^{\rm surv}\to\mathcal C_{81}\) denotes the cellular reader
of the declared enriched route, then
\[
 \widetilde{\mathfrak S}_y
 =\{x\in\Omega_{\rm HMT}^{\rm surv}:q(\pi_{81}x)=y\}
\]
is the corresponding enriched multisection. The finite certificate
proves the census and the obstruction in
\(\mathfrak S_y^{\rm atlas}\); its lift is relative to the reader
\(\pi_{81}\), not an automatic identification between a cell and a
deep history.

\begin{theorem}[Obstruction to the point selector and electronic exception]
Let \(\rho_{\rm c}(r,c')=(10-r,10-c')\) be central inversion. The thirteen fibers of
\(q\) are \(\rho_{\rm c}\)-stable. The level-zero charged-lepton fiber has
the unique fixed point, \((5,5)\). None of the other twelve fibers admits
a \(\rho_{\rm c}\)-equivariant point selector if \(\rho_{\rm c}\) acts trivially
on the coarse type.
\end{theorem}
\begin{proof}
The equation \(\rho_{\rm c}(r,c')=(r,c')\) is equivalent to \(r=c'=5\). The atlas census
places that cell in the fiber \((\ell^-,G0)\), which has cardinality one. Now let
\(y\neq(\ell^-,G0)\) and suppose that there exists an equivariant point section
\(a:\Sigma_{13}\to\mathcal C_{81}\). Since \(\rho_{\rm c} y=y\),
equivariance would require
\[
 a(y)=a(\rho_{\rm c} y)=\rho_{\rm c} a(y).
\]
Then \(a(y)\) would have to be a fixed point of \(\rho_{\rm c}\), but the unique fixed point belongs to the electron fiber. This is impossible. Therefore, in the other twelve fibers the natural output is a \(\rho_{\rm c}\)-orbit or a multisection; choosing a representative requires additional orientation or representation.
\end{proof}

The electron is exceptional in three compatible senses: \((5,5)\) is the
geometric unique center, the unique level-zero cell, and the center with respect
to which the raw lattice of differences is formed
\[
 \Lambda_{\rm raw}^{\Delta}
 =\bigl\langle A_{\rm raw}(c)-A_{\rm raw}(5,5):
 c\in\mathcal C_{81}\bigr\rangle_{\mathbb Z}.
\]
The raw signature of that cell is not null:
\[
 A_{\rm raw}(5,5)=(24,0,12,0,-12),
 \qquad \tau(5,5)=\texttt{+++---+++---}.
\]
The origin \(v_e=(0,0,0,0,0)\) of the lattice of action is a convention
relative to the electron and not the raw signature of the central cell.

The electronic algebraic invariant is evaluated by
\[
 \mathfrak e_0=
 \frac{\sqrt3}{4}
 \left(e^{\varphi/\pi^2}-22\alpha_{\rm HMT}^{3}\right)
 (1+15\Delta_4),
 \qquad
 \Delta_4=
 \frac{\pi-e\log\pi}{270}\left(1+\frac{\pi}{729}\right),
\]
\[
 \mathfrak e_0
 =0.510998922488066072509870877191349\ldots.
\]
The expression is dimensionless; the chart
\(\varsigma_{u_E}:y\mapsto y\,u_E\) subsequently introduces an energy unit.
The center \((5,5)\) and its uniqueness are exact finite results. The selection
of the coefficients \(22\) and \(15=3\cdot5\) is deduced, with explicit initial
structure, through the complement
of the five regions of \(\pi\) in \(K_e=\mathbb Z/9\mathbb Z\times\mathbb F_3\):
\(-22=-|K_e\setminus\iota(R_\pi)|\) and
\(+15=|R_\pi\times\mathbb F_3|\). The central raw signature retains a distinct type
and acts as an independent control of the basal selection.

The cell \((5,5)\) must not be confused with the energy signature
\(555555\): the latter belongs to the transversal representative of one of
the five regions of \(\pi\). They are distinct objects sharing the central
digit.

In the quark sector, the residual character
\[
 \operatorname{col}(r,c')=r-c'\pmod3
\]
refines the multisection into three classes of twelve cells and satisfies
\(\operatorname{col}(\rho_{\rm c} c)=-\operatorname{col}(c)\). This constructs the internal
color tritic. The physical representation of \(SU(3)\) is a
subsequent reader; it is not needed to prove the finite distribution.

\subsection{Stratification of extension, route, record, and signature}

The maps \(q\), \(\mathcal R\), the cellular projection, and the
realization reader distinguish four related and typed classifications:
\begin{enumerate}[label=(\roman*)]
 \item the thirteen internal family--level fibers of \(q\);
 \item the fifty-six route families of \(\mathcal R\);
 \item the eighty-one positions and their three hundred and twenty-four
       orientations;
 \item the individualized physical realizations by representation and codomain.
\end{enumerate}
The four layers are censused before the metrological contrast. For each class,
the comparison requires the tuple
\[
(\text{position},\text{fiber},\text{routes},\text{readers},
\text{record},\text{signature},\text{character},\text{towers},
\text{dimensional chart}).
\]
A reconstruction on a subset proves only that subcase and does not
replace any component of the tuple.
Given \(A\), \(C^*\), and the channels \(q_\pm\), the global tower belongs to the
infinite semigroup
\[
 s_n=q_-^n-q_+^n,
 \qquad
 n\in\langle90,120\rangle
 =\{90,120\}\cup\{30r:r\geq6\}.
\]
The heights \(90,120,180,210,240,270,360,420\) are finite representatives of the semigroup; the
preceding recurrence prolongs the complete tower. and the
dimensionless character of a signature is
\[
 \chi_{\rm mass}(v)=\exp\!\left(
 n_AA_{\rm rad}+\frac{s_C}{6}C^*_{\rm rad}
 +k\Delta_4+\nu_{120}s_{120}+\nu_{270}(135\Delta_4^2)
 \right).
\]
The moments are global: species weight them through their route, their
record, and their fibre operator. The readers \(K_D\) and \(K_\Omega\)
determine the bidirectional refinement before comparison, while the
integral table of species separately publishes the internal genealogy, the
dimensional chart, and the external comparison.

\subsection{Conjugation, anti-section, and Möbius band}

On a dodecaphasic word the exact involution is defined
\[
 C_M(\tau)_m=-\tau_{11-m},
 \qquad C_M^2=I.
\]
Linear functionals with symmetric weights are odd under \(C_M\),
whereas quadratic correlations compatible with reversion are
even. Thus, orientation and linear charge reside in the odd sector,
whereas quadratic memory resides in the even sector. This does not prove that
the complete exponential mass character is even: if
\(\chi(\tau)=e^{\ell(\tau)}\) and \(\ell\) is odd, then
\(\chi(C_M\tau)=\chi(\tau)^{-1}\).

The preceding involution acts exactly on words. To speak of an
anti-section, one additionally fixes a lift \(\widetilde C_M\) to the bundle of
sections that preserves admissibility, record, and truncations:
\[
 \operatorname{res}_{m+1,m}\widetilde C_M
 =\widetilde C_M\operatorname{res}_{m+1,m},
 \qquad
 \operatorname{sh}(\widetilde C_Mx)=C_M\operatorname{sh}(x).
\]
Only under that lift is the anti-section of \(s\)
\(\widetilde C_Ms\) on the inverted route.

The word ``Möbius'' has here a precise relative formulation. Let
\(\ell_{\rm ret}\) be a return loop of the enriched route that generates
\[
 \langle\ell_{\rm ret}\rangle\simeq\pi_1(S^1)\simeq\mathbb Z.
\]
This return group is not the finite phase \(C_9\). Suppose that its
lift to the
ledger has a sign character
\[
 \varepsilon_x:\langle\ell_{\rm ret}\rangle\longrightarrow\{\pm1\},
 \qquad \varepsilon_x(\ell_{\rm ret})=-1.
\]
The associated interval bundle is
\[
 \mathcal M_x=
 ([0,1]\times[-1,1])/(0,u)\sim(1,-u).
\]
Its first Stiefel--Whitney class is nontrivial; therefore
\(\mathcal M_x\) is a Möbius band. One circuit exchanges the two
local sheets, and two circuits restore the orientation. This does not require the
complete record to return to the same state: orientation return and memory
return remain distinct types.

The theorem is relative to the lift, the return loop, and the character
\(\varepsilon_x\), all of which are declared. It does not arise from observing two signs.
The reading of sheets additionally requires distinguishing the split grading \(R\)
from the conjugation \(C\) and imposing
\[
 CRC^{-1}=-R,
 \qquad CT(C^*)C^{-1}=T(-C^*).
\]
The particle--antiparticle reading ultimately requires the physical
representation to invert odd carriers, preserve even quadratic
carriers, and intertwine evolution and the complete character in the form declared
by the representation. Under those premises, particle and antiparticle are
the two conjugate orientations of a single persistence band.

\begin{theorem}[Even--odd decomposition on a double cover]
\label{pf84:A10-md-mobius}
Let \(p:\widetilde X\to X\) be a principal double cover with involution
\(\iota\). For every \(f\in C^0(\widetilde X,\mathbb R)\), the projectors
\[
 f^+=\frac12(f+\iota^*f),\qquad
 f^-=\frac12(f-\iota^*f)
\]
yield the decomposition \(f=f^++f^-\), with
\[
 f^+\circ\iota=f^+,\qquad f^-\circ\iota=-f^-.
\]
The even part descends to a function on \(X\); the odd part is a
section of the associated sign line. For the orientation cover
of the Möbius band, \(w_1\ne0\) obstructs a global continuous
nonvanishing section of that line.
\end{theorem}

\begin{proof}
The identities of the two projectors prove the sum and the parities. An invariant function is constant on the fibers of \(p\) and therefore descends; an anti-invariant function obeys the transition law of the sign representation. In the Möbius band, the associated line is nontrivial and its first Stiefel--Whitney class is nonzero, so every global continuous section has some zero. The conclusion excludes a global continuous section \emph{without zeros}; it does not deny the existence of sections with zeros.
\end{proof}

This classic lemma types the even--odd sector of the cover; by itself it identifies
neither word reversion, return holonomy, split
gradation, sheet conjugation, nor the parity of a
mass signature. Those identifications require their own intertwiners.

\subsection{Algebraic TRIT and temporal reader}

The declared quadratic realization associates an algebraic type with the visible sign of the TRIT. For
\(\bar\tau\in\{+1,0,-1\}\), let
\[
 \mathcal A_{\bar\tau}=
 \mathbb R[J]/(J^2+\bar\tau).
\]
Then
\[
 \begin{array}{c|c|c}
 \bar\tau&J^2&\text{regime}\ \\ \hline
 +1&-1&\text{elliptic/complex}\ \\
 0&0&\text{parabolic/nilpotent}\ \\
 -1&+1&\text{hyperbolic/split}.
 \end{array}
\]
The corresponding exponentials are
\[
 e^{tJ}=\cos t+J\sin t,
 \qquad e^{tJ}=1+tJ,
 \qquad e^{tJ}=\cosh t+J\sinh t,
\]
and in the split sector the projectors
\(P_\pm=(1\pm J)/2\) appear. These identities are verified exactly in the
declared realization. The
curvature does not arise from the isolated sign: it appears when APP supplies links and
two-cells and a connection is fixed. In the published mixed sum--product
polarization, over the \(81\) plaquettes, the oriented curvature satisfies
\[
 F(S,S)=F(P,P)=0,
 \qquad F(S,P)=+1,
 \qquad F(P,S)=-1\pmod9.
\]
The corresponding finite Stokes identity controls that the boundary of each
rectangle reproduces the sum of its plaquettes.

Ordering the transitions yields the tritic temporal reader:
\[
 \begin{aligned}
 +1&\longmapsto\text{return, memory and past},\\
  0&\longmapsto\text{torsional threshold and present},\\
 -1&\longmapsto\text{bidirectional opening and future}.
 \end{aligned}
\]
Its exact predecessor is the difference between phase return and state return:
\[
 g(k+9)=g(k),
 \qquad B_{k+9}\neq B_k\quad\text{in general}.
\]
In the corresponding geometric realization, one speaks of a spherical--volumetric sector, a flat threshold with a torsional ledger, and a hyperbolic--areal sector. This correspondence defines a Dimensional Mechanics reader subject to the map between TRIT and geometry; it neither redefines the algebraic TRIT nor by itself identifies a sign with a spacetime curvature. The expression ``time crystal'' names this architecture in HMT terminology; a dynamic realization of a time crystal additionally requires a Hamiltonian or a periodic excitation protocol, a robust subharmonic response, and stability.

\subsection{Exact support of the forward--return read}

Bidirectional exchange does not rest solely on a verbal image. The
integer angular coordinates
\[
 W_+=6n_A+s_C,
 \qquad W_-=6n_A-s_C
\]
are obtained by means of
\[
 \binom{W_+}{W_-}
 =\begin{pmatrix}6&1\\6&-1\end{pmatrix}
 \binom{n_A}{s_C},
 \qquad
 \operatorname{SNF}\!\begin{pmatrix}6&1\\6&-1\end{pmatrix}
 =\operatorname{diag}(1,12).
\]
The lattice therefore has index twelve. The involution
\(s_C\mapsto-s_C\) exchanges \(W_+\leftrightarrow W_-\) and, equivalently,
\(q_+\leftrightarrow q_-\). This is the exact algebraic substrate
of forward and return paths. Its interpretation as advanced and retarded propagation
requires the evolution operator of the physical representation.

The even projection used in the biangular reading satisfies, for every
function for which the expressions are defined,
\[
 \frac{f(\phi)+f(-\phi)}2=f_{\rm par}(\phi),
 \qquad
 \frac{(\pi-\phi)^2+(\pi+\phi)^2}{2}=\pi^2+\phi^2.
\]
They are exact symmetrization identities. Their reading in terms of a
theoretical physics model of the absorber belongs to the propagation functor, not to the
isolated algebraic identity.

\subsection{Inequivalent spinorial periodicities of \texorpdfstring{\(2\pi\)}{2 pi} and \texorpdfstring{\(4\pi\)}{4 pi}}

In the classical spinorial lift,
\[
 U(\theta)=\exp\!\left(-\frac{i\theta}{2}\sigma_z\right),
 \qquad U(2\pi)=-I,
 \qquad U(4\pi)=I.
\]
The equality is exact in the double cover \(SU(2)\to SO(3)\). Its application
to an HMT section requires specifying the spinorial representation of the
transport. Once fixed, it explains why one turn changes the sign of the
section and two turns restore it. If the CAR relations are adopted in the
occupation algebra, exclusion follows from them; neither CAR nor Pauli is
inferred from the sign of \(2\pi\) without that additional structure.

The HMT interpretation ``areal sheet \(2\pi\), volumetric sheet \(4\pi\)'' is another object. Section~\ref{sec:bisagra-hojas-pi-phi2} formalizes its core in topological terms: one circuit of the base closes the projected boundary in \(2\pi\) and ends on the conjugate sheet, whereas the squared loop closes the oriented cover in \(4\pi\). The theorem is relative to the declared lift, holonomy, and geometric reader. Its further identification with a spinorial representation still requires the physical functor and is not deduced from their sharing the same factors.

\subsection{Map between internal multisections and physical representations}

The particle object, the projection from extension to route, the internal selector of the 81 cells, the electron center, the thirteen finite family multisections, their lift relative to the cellular reader \(\pi_{81}\), the color trit, the register, the signature, the character, and the towers form the composition
\[
\text{extension}\longrightarrow\text{route}\longrightarrow
\text{multisection}\longrightarrow\text{signature}\longrightarrow
\text{character}\longrightarrow\text{towers}.
\]

The remaining obligation is to construct a realization functor that pairs certain individualized and
composite physical labels with the relevant enriched multisections and reproduces the
declared refined signatures without consulting target masses or signatures. For
composite states, the register is composed before being projected onto five
coordinates. The domain of the functor preserves the central electron, the atlas of
families, and the internal character law.


\paragraph{Evidentiary typing of persistence, atlas and families.}
Route persistence, anti-sections, the bidirectional reading, and the tritic temporal reader fix the starting architecture. The inverse system, the internal selector, and the separation of types realize that architecture through declared morphisms. The atlas censuses and the identities of \(C_M\) on words hold throughout the finite system; their lift to sections and the Möbius band are proved from explicit starting structure. The electron coefficients \((-22,+15)\) are deduced from the same explicit starting structure. If a lift first opens a metrological difference or a target signature, its value is a calibrated reconstruction.

\endgroup

\Needspace{18\baselineskip}
\section{Material ontology, charge, spin, statistics, color, and flavor}
The full definition is completed by the representation domains, the
conjugation conditions, the CAR/CCR completions, the gauge sectors,
composition, virtuality, and positivity controls.

\begingroup
\let\chapter\section
\let\section\subsection
\let\subsection\subsubsection
\let\subsubsection\paragraph
\let\clearpage\relax
\section{HMT--MD Ontology of Particles}
\Needspace{7\baselineskip}
\subsection{Material genealogical state: extension, route, record, signature, and representation}
A complete HMT state is not a bare point. It is the tuple

\[
 \mathfrak x=
 (\Gamma,\mathcal L_\Gamma,\sigma_\Gamma,\eta_\Gamma,
 \mathcal R_\Gamma,\mathcal B_\Gamma,\mathcal C_\Gamma),
\]
where \(\Gamma\) is a surviving route, \(\mathcal L_\Gamma\) its genealogical record, \(\sigma_\Gamma\) its integer signature, \(\eta_\Gamma\) its harmonic refinement, \(\mathcal R_\Gamma\) its representation, \(\mathcal B_\Gamma\) its return memory, and \(\mathcal C_\Gamma\) its composition or boundary datum. Two states with the same visible word may be distinct if they differ in prefix, block, memory, carry, sheet, or Witt shadow.

\Needspace{7\baselineskip}
\subsection{Particle as a persistent local section of the dimensional genealogical state}
A particle is an equivalence class of persistent sections that satisfies four conditions:

\begin{enumerate}
\item compatible prolongation through the TPK cylinders;
\item phase return with state memory;
\item irreducible physical representation or minimal invariant block;
\item spectral stability under the coupling that observes it.
\end{enumerate}
A persistent section is a compatible family.

\[
 s=(s_K)_{K\ge K_0},\qquad
 p_{L,K}(s_L)=s_K\quad(L\ge K\ge K_0).
\]
Two sections \(s\) and \(s'\) representan the same particle when there exists a level \(K_1\) and a family compatible of simetrias internal \(g_K\in G_K\) tal that

\[
 s'_K=g_Ks_K\quad(K\ge K_1),
 \qquad
 p_{L,K}g_L=g_Kp_{L,K},
\]
and the observable characters of the genealogical record coincide. This relation identifies descriptions that differ by internal gauge, but it does not identify routes with distinct memory, carry, sheet, block, or holonomy. The particle is the class

\[
 [s]_{\rm pers}
 =\{s':s'\sim s\},
\]
not a single visible word nor a cell chosen after knowing its mass.

If \(\widetilde\Gamma_K\) is the extension at level \(K\), persistence is expressed by

\[
 p_{K+9,K}(\widetilde\Gamma_{K+9})
 =\widetilde\Gamma_K,\qquad
 g(K+9)=g(K),
\]
without requiring

\[
 \mathcal B_{K+9}=\mathcal B_K.
\]
The particle is not the returning phase, but the genealogical identity preserved through the change of state. Its mass measures the spectral rigidity of that persistence when it is coupled to a material realization.

\Needspace{7\baselineskip}
\subsection{Antiparticle as the conjugate image of a persistent section under sheet inversion}
The dodecaphonic conjugation is

\[
 C_M(\mathbf t)=-\operatorname{rev}(\mathbf t).
\]
The particle--antiparticle pair is formed by two orientations of the same genealogical band. If \(\chi\) denotes the sheet orientation,

\[
 (\mathbf t,\chi)\sim(C_M\mathbf t,-\chi).
\]
For a band energy

\[
 E_{\rm band}(\mathbf t,\chi)
 =E_+(\mathbf t)+\chi E_-(\mathbf t),
\]
with

\[
 E_\pm(\mathbf t)
 =\frac12\bigl(E_0(\mathbf t)\pm E_0(C_M\mathbf t)\bigr),
\]
we obtain

\[
 E_{\rm band}(C_M\mathbf t,-\chi)
 =E_{\rm band}(\mathbf t,\chi).
\]
The equality of masses is therefore a consequence of the operator's parity under conjugation, while charge and orientation change. Equality of widths additionally requires the decay operator to be conjugated by \(C_M\).

\Needspace{7\baselineskip}
\subsection{Mass as a spectral character of persistence of \(1\) a material route}
Mass is the spectral measure, on the line \(\mathcal L_M\), of the operator resulting from compressing the character, towers, bidirectional reader, and clock onto the persistent subspace selected by the coupling:

\[
 \widehat M_{P,a}
 =P_{P,a}\,
 \mathbf m_\star^{\rm ret}
 R_{\rm act}^{\widehat K_{P,a}/2}
 \widehat{\mathcal A}_{P,a}^{(0)}
 R_{\rm act}^{\widehat K_{P,a}/2}
 P_{P,a}.
\]
The general output is

\[
 \mu_{P,a}^{\rm mass}(B)
 =\operatorname{Tr}\!\left(
 \rho_{P,a}E_{\widehat M_{P,a}}(B)\right).
\]
A scalar mass \(m_P\) exists if and only if

\[
 P_{\operatorname{supp}\rho}\widehat M_{P,a}
 P_{\operatorname{supp}\rho}
 =m_P P_{\operatorname{supp}\rho},
\]
equivalently, if the variance of \(\widehat M_{P,a}\) in the prepared state is zero. This definition includes the rank-one electron, the multiplets forced by Schur's lemma, and isolated resonance poles, without conflating them.

\Needspace{7\baselineskip}
\subsection{Charge as a transported caliber character by the material section}
Charge is the oriented character of the odd layer of the genealogical record under conjugation. Let \(\sigma_{\rm odd}\) be the odd signature. Then

\[
 Q(C_M\Gamma)=-Q(\Gamma),\qquad
 Q(\Gamma)=\mathfrak q\!\left(\sigma_{\rm odd}(\Gamma)\right),
\]
where \(\mathfrak q\) is the homomorphism from the odd signature to the charge lattice of the realization. The even mass signature and the odd charge signature are distinct objects: a word may have zero linear sum and nevertheless carry a nonzero mass signature.

\Needspace{7\baselineskip}
\subsection{Spin as a representation of the rotational action on the material fiber}
Spin is the holonomy class with which a persistent section responds to a complete rotation of the orientation frame. The internal root

\[
 J^2=-I
\]
separate the two closures:

\[
 U(2\pi)=
 \begin{cases}
 +I,&\text{tensorial sector},\\
 -I,&\text{spinorial sector},
 \end{cases}
 \qquad
 U(4\pi)=I.
\]
An integer-spin boson lies in a representation that closes at \(2\pi\); a half-integer-spin fermion requires the double circuit \(4\pi\). The spin value is not deduced from an isolated cardinal count: it is the weight of the representation of the holonomy group on the persistent subspace. In particular:

\[
 s=0:\text{ scalar section},\qquad
 s=\tfrac12:\text{ spinorial band},
\]
\[
 s=1:\text{ transverse vector section},\qquad
 s=2:\text{ tensor component of a metric realization}.
\]
The last row defines a representation type; it does not require the existence of an elementary spin-two particle.

Chirality and helicity are subsequent and distinct data. An oriented involution \(\Gamma_{\rm ch}^2=I\) defines

\[
 P_L=\frac{I-\Gamma_{\rm ch}}2,
 \qquad
 P_R=\frac{I+\Gamma_{\rm ch}}2,
\]
whereas helicity is the spectral sign of the rotation generator along the direction of propagation. The former classifies sheets of the representation; the latter depends on the kinematic state. This separation allows a neutrino to be fermionic through its closure \(4\pi\), even though its weak coupling selects a single chirality.

\Needspace{7\baselineskip}
\subsection{Exchange symmetry induced quantum statistics}
Statistics is the rule of completion of the space of identical paths:

\[
 \mathcal F_-(\mathcal H)
 =\bigoplus_{n\ge0}\wedge^n\mathcal H,\qquad
 \mathcal F_+(\mathcal H)
 =\bigoplus_{n\ge0}\operatorname{Sym}^n\mathcal H.
\]
In the exterior completion, the CAR relations

\[
 \{a_i,a_j^\dagger\}=\delta_{ij},\qquad
 \{a_i,a_j\}=0
\]
imply

\[
 n_i^2=n_i
\]
and therefore Pauli exclusion and the Fermi--Dirac distribution. In the symmetric completion, the CCR relations permit multiple occupancy and lead to Bose--Einstein. Statistics are not inferred from mass: they follow from how the route is composed under exchange.

\Needspace{7\baselineskip}
\subsection{Color as residual representation of \(SU(3)\) on the quark routes}
Color is the residual triadic orbit of a quark section within the charge and flavor fiber. The thirty-six colored positions decompose as

\[
 12_R+12_G+12_B.
\]
The physical object is not one of those colors, but the triple and its action:

\[
 \mathcal H_q^{\rm color}
 =\mathcal H_R\oplus\mathcal H_G\oplus\mathcal H_B.
\]
If the mass operator commutes with the irreducible color action,

\[
 U_g\widehat M_qU_g^*=\widehat M_q,
\]
Schur's lemma forces a common mass in the triplet. Choosing a red, green, or blue cell as representative would destroy covariance.

\Needspace{7\baselineskip}
\subsection{Taste as a class of route within a material generation}
Flavor is an inequivalent coupling class within a single gauge representation, characterized by charge, generation, leaf, fine signature, memory, and interaction projector. Two flavors may share a coarse fiber and differ in the intertwiner coupling them to the reader.

Formally, if \(F\) is the common fiber, the flavors \(f_i\) are representations \(\pi_i\) for which

\[
 \operatorname{Hom}_{G_F}(V_{f_i},\ell^2(F))\ne0
\]
and whose projectors \(P_{f_i}\) are not conjugate under a gauge symmetry that remains physical. Flavour is not the depth of the atlas and does not reduce to a conventional name.

\Needspace{7\baselineskip}
\subsection{Material generation as a typed recurrent layer, stratified by return depth, fine signature, and reader spectrum}
A generation is a recurrent stratum of the same grammar of charge and representation, distinguished by return depth, fine signature, and reader spectrum. The generation label is physical only when there is a projector that intertwines that stratum with the coupling channel. A depth \(G_j\) must therefore not be identified automatically with the electron, muon, tau, or sterile neutrino.

\Needspace{7\baselineskip}
\subsection{Gauge field as a connection on the charge fibers}
A gauge field is the realization of a frame redundancy over the fibers. Its elementary bosons belong to the null direction of the mass operator while the symmetry remains exact. A gauge mass is not obtained by approximating a positive character to zero, but through a breaking, a longitudinal coupling, or a declared pole realization.

\Needspace{7\baselineskip}
\subsection{Composite particle as a persistent section of the product of linked paths}
A composite particle is a persistent section of the composition operator, not the sum of the masses or signatures of its constituents. If \(\Gamma_1,\ldots,\Gamma_n\) are constituent routes, first construct

\[
 \mathfrak C_{12}
 (\Gamma_1,\ldots,\Gamma_n;\mathcal T,\partial,\mathrm{carry})
 =\mathcal L_{\rm comp},
\]
where \(\mathcal T\) is the temporal tree and \(\partial\) the link boundary. Only afterwards does the reader act

\[
 \mathcal L_{\rm comp}\longrightarrow\sigma_{\rm comp}
 \longrightarrow\chi_{\rm comp}
 \longrightarrow\widehat M_{\rm comp}.
\]
The binding energy resides in the composition cocycle, boundary memory, and modification of the spectrum; it is not introduced as an external correction.

\Needspace{7\baselineskip}
\subsection{Resonance as a transient section and width as a decay rate}
A resonance is a persistent non-unitary mode of the coupled operator. If

\[
 \lambda=re^{-i\theta},\qquad 0<r<1,
\]
then

\[
 E=\frac{\hbar_{\rm int}}{t_0}
 \bigl(\theta-i\log r\bigr)
 =M-\frac{i}{2}\Gamma.
\]
The mass \(M\) and width \(\Gamma\) are the real and imaginary parts of a pole, but they require distinct readers. The width may also be obtained from a compressed semigroup:

\[
 T_X=Q_XUQ_X,\qquad
 \lambda_X=-\frac1{t_0}\log\rho(T_X),\qquad
 \Gamma_X=\hbar_{\rm int}\lambda_X.
\]
No width is deduced solely from the fact that a route possesses finite horizon.

\Needspace{7\baselineskip}
\subsection{Virtual section as an intermediate non-asymptotic extension of the atlas}
A section is virtual when its prolongation ends in a non-trivial manner:

\[
 \operatorname{Ext}_{L}(s)\ne\varnothing,\qquad
 \operatorname{Ext}_{L+1}(s)=\varnothing,\qquad
 \mathcal L(s)=L<\infty.
\]
Its primary content is the horizon \(L\), the terminal obstruction, the finite genealogical register, and the propagation amplitude. It is neither an imaginary-mass particle nor an as-yet-unclassified documentary state.

\Needspace{7\baselineskip}

\endgroup
\begingroup
\let\chapter\section
\let\section\subsection
\let\subsection\subsubsection
\let\subsubsection\paragraph
\let\clearpage\relax
\section{Complete classification of particles and physical states}
\label{sec:catalogo-realizaciones-tipadas-rev8}
\index{physical sectors!typed inventory}
\index{mass atlas!81/56/13/324}

The realizations are ordered by object, multiplicity, and codomain. The mathematical domain contains \(81\) positions, grouped into
distinct codomains. The mathematical domain contains \(81\) positions, grouped into
\(56\) route families and \(13\) extension classes; the \(324\)
oriented routes adjoin the initial orientation and do not represent additional masses
or particles. The ``evaluation operator'' column does not impose a mass where the object requires
a null chart, a prior composition, a topological invariant, or a
prolongation obstruction.

\subsection{Spectral definitions of particle, spin, color, and flavor}

A \emph{particle} is a surviving section realized as an oriented route, endowed with a register, a symmetry representation, a spectral operator, and a physical chart. A fermion realizes the exterior grading and the CAR relations; a boson realizes the symmetric completion and the CCR relations. Spin records the holonomy of orientation: the return of \(2\pi\) closes the bosonic phase, whereas the spinorial section recovers its state after \(4\pi\). This distinction links the double areal--volumetric sheet to parity and Pauli exclusion.

The \emph{color} is the ternary representation that lifts the coordinate
\(r-c\pmod3\) to the fundamental orbit and thereafter to the adjoint of
\(\mathfrak{su}_3\). The \emph{flavor} identifies a realization basis
within a representation with fixed charge, color, and spin; CKM and PMNS
transport between the interaction basis and the spectral basis. Neutrinos
are neutral fermionic sections: their flavor labels belong to the
mixing reader and their masses to the eigenvectors of the corresponding operator.

A virtual particle is a finite section with a ledger and propagator up to
a horizon of prolongation; a Majorana realization is the real
fermionic sector fixed by conjugation; and the Fibonacci, Ising, and
\(\mathbb Z/6\mathbb Z\) sectors publish quantum dimensions, fusion rules, and
braid characters. Each definition thus preserves the codomain assigned to it by the
structure.


\endgroup
\begingroup
\let\chapter\section
\let\section\subsection
\let\subsection\subsubsection
\let\subsubsection\paragraph
\let\clearpage\relax

The table also preserves the quantum realizations abbreviated by its
classification. The three charged leptons and the six quark flavours belong to the
exterior completion: CAR leads to Pauli and Fermi--Dirac, without attributing those
relations to an isolated TRIT. The quarks occupy thirty-six cells when
colour and antiparticle are conserved before the quotient. The common neutral support
\(\mathcal H_\nu=\bigoplus_{j=1}^{4}\mathcal H_{\nu,G_j}\) has dimensions
\(2+4+6+8=20\). The symbols \(G_j\) designate depth strata of the atlas;
\(f_e,f_\mu,f_\tau\) designate the interaction basis. In particular, \(G_4\) is
a stratum of the common neutral support and is not identified with a sterile or
exotic state without an additional projector. PMNS transports the interaction basis to the
eigenvector basis; it creates neither eigenvalues, a mass scale, nor a
depth selector. The photon preserves outward passage and return within the
null gauge chart: \(m=0\) does not mean evaluating \(\chi_0\) with a highly
negative exponent. \(W^\pm\) and \(Z^0\) are massive bosonic representations, with
charge conjugation for \(W^+\leftrightarrow W^-\). Lastly, the
antisection is obtained through inversion and band reversal; mass is even
under that involution and is not duplicated as a second independent parameter.

All sectors preserve the causal prefix.
\[
\text{surviving extension}\longrightarrow\text{observable route}
\longrightarrow\text{enriched chronological record}.
\]
The reader then changes. In the positive material sector, it continues through
\(\text{signature}\to\text{character}\to\text{towers}\to\text{dimensional chart}\).
The null sector bypasses the positive character and fixes \(m=0\); hadrons first compose
enriched chronological ledgers; topological sectors use fusion or
braid invariants; and virtuality is recognized by terminal obstruction. The
domain of these maps is the complete preceding classification. A
numerical trial on a subset establishes only the finite case tested and does not
reduce that domain.

The censuses accompanying this table belong to different quotients:
\[
81=25_{\ell^-}+20_\nu+20_d+16_u,
\]
\[
56\ \text{families with fibers}\quad
41\times1+10\times2+2\times3+1\times4+2\times5=81,
\]
\[
13\ \text{multisections},\qquad324=81\times4\ \text{genealogies}.
\]
Likewise, the \(192\)-A representation classes associated with
\(\operatorname{diag}(2,6,2,2,4)\) and the lattice index \(192\)-B of
\(\operatorname{SNF}(2,2,2,2,12)\), with \(52\) raw signatures, are not the
The cardinality of the historical expository catalogue does not count individual particles or define the material taxonomy.
not individual particles.

\endgroup

\begingroup
\let\chapter\section
\let\section\subsection
\let\subsection\subsubsection
\let\subsubsection\paragraph
\let\clearpage\relax
% Vista científica exacta materializada desde una rama condicional.
% Fuente: source/demostraciones_primarias/27_04_particula_estadistica.tex; rama: HMTIncludeParticleCore.
% No modifica el contenido matemático de la rama de origen.
\chapter{Persistent Sections and Finite Propagation}
\label{chap:particula}

The mathematical notion of particle is introduced after the boundary structure and before its dynamics are specialized. Its primary support is a surviving extension of the TPK\@ state system. An oriented trajectory is the sequential shadow of that extension; over it, the internal degrees form a discrete bundle. The propagator, the sum over paths, and the Clifford representations are incorporated after the line of action has been constructed. This hierarchy prevents a particle from being defined through the external selection of a row or route and separates the mathematical object from its subsequent identification with a physical species.

\section{Surviving extension and observable projection}

Let \((\mathcal S_m,\rho_{m+1,m})\) be the inverse system of states that
pass the closures up to depth \(m\). A global state is a compatible
family
\[
 x=(x_m)_m\in\Omega_{\HMT}^{\rm surv}
 :=\varprojlim_m\mathcal S_m.
\]
A sequential chart of depth \(R\) produces a route
\[
 \gamma_x^{[R]}=\operatorname{sh}_R(x)
 =(c_0\longrightarrow\cdots\longrightarrow c_R).
\]
The symbol \(\operatorname{sh}\) denotes a shadow reading: different charts may
present the same global state, provided that they preserve the sheet,
orientation, boundary, and record. In \cref{chap:extension-ruta-caracter}, the enriched record is
constructed so that the action character descends from these presentations
to the extension.

\section{Particle as a compatible section of a discrete bundle}

Let \(\Gamma_{\rm TPK}\) be the graph whose vertices are complete states of the
Topological Phase Kernel and whose edges are admissible transformations.
For \(x\in\Omega_{\HMT}^{\rm surv}\) and its observed trajectory
\[
 \gamma_x=(c_0\longrightarrow c_1\longrightarrow\cdots
 \longrightarrow c_R),
\]
let \(\operatorname{Typ}(c_t)\) be the finite set of arithmetic types,
orientations, and memories over the state \(c_t\), and let
\(E_{c_t}=\mathbb C[\operatorname{Typ}(c_t)]\) be the free complex vector space
spanned by those types. The discrete bundle restricted to \(\gamma_x\) is
\[
 \mathcal B_{\gamma_x}=\bigsqcup_{t=0}^{R}E_{c_t}
 \longrightarrow\{0,\ldots,R\}.
\]
Each edge \(c_t\to c_{t+1}\) induces a linear transport map
\(\tau_t:E_{c_t}\to E_{c_{t+1}}\).

\begin{definition}[Mathematical particle]
A particle of Dimensional Mechanics is a tuple
\[
 \mathfrak p=(x,\gamma_x,\mathcal B_{\gamma_x},s,n,U),
\]
where \(x\in\Omega_{\HMT}^{\rm surv}\) is the surviving extension and
\(\gamma_x=\operatorname{sh}_R(x)\) its shadow in the stated chart; moreover,
\[
 s:\{0,\ldots,R\}\longrightarrow\mathcal B_{\gamma_x},\qquad
 s(t)\in E_{c_t},\qquad s(t+1)=\tau_t(s(t)),
\]
is a compatible section;
\[
 n:\bigsqcup_{t=0}^{R}\operatorname{Modes}(E_{c_t})
 \longrightarrow\mathbb N
\]
is its occupation function, where
\(\operatorname{Modes}(E_{c_t}):=\operatorname{Typ}(c_t)\) denotes the distinguished
basis of the free vector space; and
\[
 U_{t_2,t_1}:E_{c_{t_1}}\longrightarrow E_{c_{t_2}}
\]
is a family of linear propagators that intertwines the transports,
\(U_{t,t}=\operatorname{id}_{E_{c_t}}\),
\(U_{t+1,t}=\tau_t\), and satisfies
\[
 U_{t_3,t_2}\circ U_{t_2,t_1}=U_{t_3,t_1}
 \qquad(0\leq t_1\leq t_2\leq t_3\leq R).
\]
Each composition is defined between the corresponding fibers.
The identification of \(\mathfrak p\) with a physical
species additionally requires a representation of the pertinent symmetry group
and a map from HMT states to the states of that representation.
\end{definition}

The extension preserves the genealogy and the register; the route exposes its ordered reading; the internal degrees belong to the fibers, occupation to \(n\), and evolution to \(U\). The particle is the compatibility of these levels, not any one of them separately. In particular, the route is not chosen by the particle: it is a projection of the extension that realizes it.

\endgroup
\begingroup
\let\chapter\section
\let\section\subsection
\let\subsection\subsubsection
\let\subsubsection\paragraph
\let\clearpage\relax
% Fragmento científico exclusivo de 01_electron.tex, líneas 120--197.
% El prefijo 1--118 coincide byte a byte con la ontología principal ya
% publicada; este fragmento conserva únicamente la formulación cuántica
% adicional sin repetir definiciones idénticas.
\section{Quantum Ontology of the Particle}
\Needspace{7\baselineskip}
\subsection{State, history and realization}
A genealogical state is a compatible section of the TPK system of extensions. An observable history is the chronological projection of that section onto a reader. A physical realization is the representation of the history on a dimensional line, a Hilbert space, a field algebra, or a topological codomain. These three notions are not interchangeable. The state carries identity; the history carries evidence; the realization carries the observable.

Let \(\mathcal X_{\rm TPK}^{\rm enr}\) be the space of enriched states and \(p\) the observable reader. A particle is a class \([X]\) of persistent sections such that

\[
p(X_{K+1})|_K=p(X_K),
\]
the transition law preserves the relations of the genealogical record, and the prolongation does not terminate at an obstruction horizon. The route

\[
\gamma_X=(p(X_0),p(X_1),p(X_2),\ldots)
\]
is the history visible of the particle. Not is a trajectory espacial by definition. A realization metric posterior may convertirla in curve, orbit or propagador when constructs the entrelazador corresponding.

The physical identity is preserved under refinement when the squares

\[
\begin{CD}
X_{K+1} @>{F_{K+1}}>> X'_{K+1}\\
@V{r_{K+1,K}}VV @VV{r'_{K+1,K}}V\\
X_K @>>{F_K}> X'_K
\end{CD}
\]
commute. Mass, charge, and spin are different readers of that same antecedent. Equality of one does not entail equality of the others.

\Needspace{7\baselineskip}
\subsection{Quantum representation of the particle as a persistent local section of the genealogical state}
An HMT--MD particle is a persistent section of an enriched extension satisfying four conditions:

\begin{enumerate}
\item has an observable route compatible with the nonadic monodromy;
\item preserves a genealogical record under refinement;
\item admits a representation of orientation, charge, spin, and statistics;
\item presents a stable spectral output in the declared physical codomain.
\end{enumerate}
Stability does not mean immobility. It means that changes of phase, sheet,
and memory belong to the same genealogical class. The particle may close
its phase but not its state, change sheets while preserving its mass, or
traverse several cells while retaining the same extension identity.

A species is a family of particles that shares the internal classifier.

\[
\mathfrak s=
(\mathfrak q,\mathfrak j,\mathfrak c,\mathfrak f,\mathfrak g,
\mathfrak r,\mathfrak o),
\]
where \(\mathfrak q\) is the charge representation, \(\mathfrak j\) the spin holonomy, \(\mathfrak c\) the color orbit, \(\mathfrak f\) the flavor, \(\mathfrak g\) the generation, \(\mathfrak r\) the route family, and \(\mathfrak o\) the mass operator. A non-central species may correspond to a complete multisection; requiring it to have a single APP point before declaring the physical coupling is a type error.

\Needspace{7\baselineskip}
\subsection{Quantum conjugation of the antiparticle via sheet inversion of a persistent section}
The dodecaphasic involution is

\[
C_M(\tau)=-\operatorname{rev}(\tau).
\]
It acts on the oriented history and prolongs to the genealogical ledger. Charge and linear orientation are odd; quadratic memory is even. If \(q\) is the charge reader,

\[
q(C_MX)=-q(X).
\]
Equality of masses further requires operatorial covariance. If \(K\) realizes the conjugation in the physical space and

\[
K\widehat M_XK^{-1}=\widehat M_{\bar X},
\]
also transporting the channels and the causal prescription when \(K\) is antiunitary, then \(m_{\bar X}=m_X\). The parity of the record by itself does not replace this premise.

Particle and antiparticle are two oriented sections of the same genealogical band. In the Möbius realization, one circuit interchanges the sheet and two circuits restore the orientation. The conclusion does not follow from drawing a twisted ribbon: it requires the involution, the lift to the double cover, the return loop, and the sign holonomy.

For the electronic record \(\tau_e=+++---+++---\),

\[
C_M\tau_e=\tau_e.
\]
The central word is fixed, while the charge representation distinguishes the electron from the positron. This yields the same mass without duplicating the temporal route. Closure after 216 steps does not mean 108 particle steps and 108 antiparticle steps: it means restoration of the oriented kinematic lift.

\Needspace{7\baselineskip}

\endgroup
\begingroup
\let\chapter\section
\let\section\subsection
\let\subsection\subsubsection
\let\subsubsection\paragraph
\let\clearpage\relax
The General Mass Law is formulated as a genealogical and
spectral transformation, not as a collection of numbers. Its effective domain is the chain
\[
\text{APP}\longrightarrow\text{TRIT}\longrightarrow\text{TPK}
\longrightarrow\Gamma^{\rm surv}\longrightarrow\Gamma^{\rm obs}
\longrightarrow L_{108}\longrightarrow\sigma\in\mathbb Z^5
\longrightarrow\chi_{\rm full}\longrightarrow\widehat M,
\]
followed, when the physical representation permits it, by the functional that
extracts a scalar coordinate from the mass line. Each arrow preserves its
domain, orientation, memory, and status; no external figure retroactively selects
a route, signature, or harmonic coefficient.

The finite atlas contains 81 cells, 56 route families, 13 multisections, and
324 oriented routes. The central electronic point is the unique equivariant point
fiber. The remaining elementary classes are naturally realized as
fibers, orbits, or multisections, so their primary output may be an
operator or spectrum before being a scalar mass. This difference is not a
deficiency of domain, but a property of the internal geometry.

In the central fiber, the basal selector and the bidirectional reader produce
\[
\mathfrak e_0
=\frac{\sqrt3}{4}
 \left[\exp\!\left(\frac{\varphi}{\pi^2}\right)
       -22\alpha_{\rm HMT}^3\right](1+15\Delta_4),
\qquad
\kappa_e=80-54\alpha_{\rm HMT}+6\Delta_4,
\]
\[
\mathfrak e_e^\beta=\mathfrak e_0R_{\rm act}^{\kappa_e},
\qquad
\varsigma_{u_E}(\mathfrak e_e^\beta)
=0.510998950690000808608\ldots\ \mathrm{MeV}.
\]
The absolute action scale remains on the
\(\Sigma_H=\varphi\alpha_{\rm HMT}^{16}\) branch, of decimal order \(-34\), and is
coupled to the clock through \(\hbar_{\rm int}\omega\); when mass ratios are formed,
every common section cancels and only the correctly typed characters,
readers, and relative refinements survive.

The harmonic hierarchy does not consist of eight isolated corrections. It is the global semigroup.
\[
\langle90,120\rangle
=\{90,120\}\cup\{30r:r\ge6\},
\qquad
c_n=q_-^n+q_+^n,
\quad s_n=q_-^n-q_+^n,
\]
acting on a single bilayer operator. The visible levels are witnesses to
that hierarchy; their representational completeness does not replace the physical map that
must transport route, carry, memory, and winding to the coefficients of each
state.

The typed physical realizations, the seventeen conditioned evaluations, the two null sections, and the neutrino, composite, resonant, topological, and virtual sectors are thus brought together without flattening their codomains. The subsequent inventory of 471 masses and 384 widths preserves the state-by-state comparison but is not conflated with a census of predictions. A quantitative comparison acquires physical force only after its observable, unit, scheme, scale, uncertainty, covariance, and realization rule have been declared.

The structural conclusion is precise: mass is the spectral measure of persistence under coupling. The particle is the persistent section; the route, its observable history; the signature, its integer reading; the character, its dimensionless action; the towers, its harmonic resolution; and the dimensional chart, the realization of that genealogy in physical units. The preceding atlas and inventories permit each claim to be reproduced without reducing the theory to an abbreviated table or turning an open interface into a closed theorem.

\endgroup

\section{Composition and identity}
A composite particle retains its identity when the binding record, the
boundary, the representation, and the observables remain compatible under
refinement. A resonance adds a pole and a width; it does not constitute a
new coordinate of the mass signature. The identity criterion is refuted
if it depends on the external value or if two inequivalent histories are identified
without an isomorphism of sections.
